% Guard: macros.tex is \input both by main.tex (after the packages it needs,
% e.g. amsmath for \DeclareMathOperator) and by individual section files so
% they remain compilable on their own. Loading it twice would redefine every
% \newcommand and abort, so the second and later loads are no-ops.
\ifdefined\PQRmacrosloaded\endinput\fi
\let\PQRmacrosloaded\relax

% --- 0. revision marking ---
% Text added to Section 2 during the current revision is wrapped in \added{...}
% so the advisors can see at a glance what is new. Redefine to \relax to print
% the paper without the marking.
\newcommand{\added}[1]{\textcolor{blue}{#1}}
\long\def\added#1{\begingroup\color{blue}#1\endgroup}  % \long: may span paragraphs


% --- A. operator-style score and function names ---
\newcommand{\HD}{\mathit{HD}}                        % histogram distance; use \HD_\gamma
\newcommand{\Attr}{\mathit{Attr}}                    % attribute set of a schema
\newcommand{\Attrs}{\mathit{Attrs}}                  % attribute set argument
% Abbreviated in MATH only; the prose still spells the names out in full.
% Spelling them out inside formulas makes displayed equations far too wide
% for a two-column layout (the product-rule step below overflowed by 200pt).
\newcommand{\Use}{\mathit{Use}}
\newcommand{\Cov}{\mathit{Cov}}
\newcommand{\Rep}{\mathit{Rep}}
\newcommand{\DivS}{\mathit{Div}}
\newcommand{\TVD}{\mathit{TVD}}
\newcommand{\GS}{\mathit{GS}}                        % global sensitivity
\DeclareMathOperator{\Dom}{Dom}
\DeclareMathOperator{\Lap}{Laplace}

% --- B. core objects ---
\newcommand{\Qone}{Q_1}
\newcommand{\Qtwo}{Q_2}
\newcommand{\pD}{p(D)}
\newcommand{\views}{\mathcal{P}}
\newcommand{\hist}[2]{H_{#1,#2}}                     % \hist{\Qone(\pD)}{A}
\newcommand{\dhist}{\Delta H}
\newcommand{\nhist}{\widetilde{\Delta H}}            % the released noisy object

% --- C. privacy budget and parameters ---
\newcommand{\eps}{\epsilon}
\newcommand{\epsP}{\epsilon_P}
\newcommand{\epsS}{\epsilon_S}
\newcommand{\epsH}{\epsilon_H}
\newcommand{\gammaref}{\gamma_{\mathrm{ref}}}
\newcommand{\gammaeval}{\gamma_{\mathrm{eval}}}
\newcommand{\taumax}{\tau_{\max}}

% --- D. mechanisms ---
\newcommand{\ExpMech}{\textsc{ExpMech}}
\newcommand{\OneShotTop}[1]{\textsc{OneShotTop-}#1}
\newcommand{\AdaptiveTau}{\textsc{AdaptiveTau}}
% --- E. experiment methods (used in sections/experiments.tex); a method is renamed here ---
% (the figure legends are drawn by src/pqde/paper_figures.py and must be redrawn after a rename)
\newcommand{\PrivExpl}{\textsf{PrivExpl}}
\newcommand{\OptPrivExpl}{\textsf{OptPrivExpl}}
\newcommand{\BF}{\textsf{BF}}
\newcommand{\SingleStage}{\textsf{SingleStage}}
\newcommand{\RandomExpl}{\textsf{Random}}
\newcommand{\Bone}{\ensuremath{\mathsf{B1}}}
\newcommand{\Btwo}{\ensuremath{\mathsf{B2}}}
\newcommand{\Bthree}{\ensuremath{\mathsf{B3}}}

\section{Explanations and Their Quality}\label{sec:model}

We now define what an explanation of a query difference is (Section~\ref{subsec:private-query-difference-setting}) and what makes one good (Section~\ref{sec:explanation-quality}).
Section~\ref{sec:solution} then shows how good explanations are selected and released under differential privacy.

\subsection{Private Explanations of Query Differences}
\label{subsec:private-query-difference-setting}

\paratitle{Setting.}
A trusted server holds a private database $D$ and receives two related queries $Q_1$ and $Q_2$ over it.
The analyst wants to know \emph{where in the data the two outputs differ, and in what respect}.
In the running example (Example~\ref{ex:intro}), the two outputs share only $874$ of their tuples, and the analyst wants to know which populations the refinement removed, which it added, and how it reshaped those it kept.
Since $D$ is private, the answer has to be released under differential privacy.

\paratitle{Query refinement.}
We focus on pairs in which $Q_2$ is a \emph{refinement} of $Q_1$.
In such a pair, the output of $Q_1$ violates constraints on its composition, such as bounds on the number or the fraction of output tuples that belong to a group.
The query $Q_2$ is obtained by relaxing, tightening, removing or replacing atomic conditions of $Q_1$ so that its output satisfies these constraints, typically with a minimal change to the predicate~\cite{MK09,LMSJ23,ERICA,CSSM24}.
Example~\ref{ex:intro} is of this kind.
We treat the refinement mechanism as a black box.
The model uses only the fact that both queries are selections over $D$, so it applies equally to two groups of a table that the analyst wishes to compare~\cite{BayPazzani01,TaoGMR22} (Section~\ref{subsec:ext-contrast}) and to two versions of a query written by hand.
Section~\ref{subsec:ext-cqa} discusses outputs that are not selections.
%For a view $p$, we write $Q_i(\pD)$ for the output of $Q_i$ restricted to the view, which equals $Q_i(D)\cap\pD$ (Eq.~\eqref{eq:view-restriction}).

\paratitle{Workflow.}
Figure~\ref{fig:framework} shows the interaction.
The analyst submits the two queries, a collection $\views$ of candidate view predicates, the explanation attributes $\Attrs$, the reference attributes $\mathcal{A}_R$, the number $k$ of explanations, and a privacy budget $\epsilon$.
The server evaluates both queries, privately selects $k$ explanations, and releases them.
Nothing else leaves the server.
% Both queries are inputs, so 
If the mechanism that produced $Q_2$ read $D$, its privacy cost composes with ours and is outside our scope.

\begin{figure}[t]
\centering
% Architecture of PEEQ (Figure 1). Single-column TikZ replacement for
% figures/framework.png. Requires tikz with the libraries positioning,
% arrows.meta, fit and backgrounds (loaded in main.tex).
\definecolor{peeqblue}{RGB}{45,95,150}
\definecolor{peeqorange}{RGB}{196,88,32}
\begin{tikzpicture}[
  font=\scriptsize,
  >={Stealth[length=1.5mm,width=1.2mm]},
  stage/.style={draw=peeqblue, line width=0.5pt, rounded corners=1.5pt, fill=white,
                align=center, inner xsep=1.5pt, inner ysep=2.5pt, minimum height=11mm},
  flow/.style={->, line width=0.5pt, draw=black!75},
]
% --- the three private steps inside the server ---
\node[stage, text width=19mm] (view)
  {\textbf{View stage}\\ top-$\tau$ views\\ by $\Use$ \,{\color{peeqblue}($\epsilon_P$)}};
\node[stage, text width=26mm, right=3mm of view] (expl)
  {\textbf{Explanation stage}\\ $k$ rounds by $\HD_\gamma\cdot\DivS$\\ {\color{peeqblue}($\epsilon_S$)}};
\node[stage, text width=19mm, right=3mm of expl, draw=peeqorange, fill=peeqorange!10] (rel)
  {\textbf{Release}\\ noisy pairs of\\ histograms \,{\color{peeqorange}($\epsilon_H$)}};
\draw[flow] (view) -- (expl);
\draw[flow] (expl) -- (rel);
% --- trusted server ---
\node[anchor=south west, inner sep=0pt, text=peeqblue] (srvlabel)
  at ([yshift=1.8mm]view.north west)
  {\textbf{Trusted server} (holds $D$),\; $\epsilon=\epsilon_P+\epsilon_S+\epsilon_H$};
\begin{scope}[on background layer]
\node[draw=peeqblue!70, dashed, line width=0.5pt, rounded corners=3pt, fill=peeqblue!6,
      inner xsep=1.5mm, inner ysep=1.5mm, fit=(view)(expl)(rel)(srvlabel)] (server) {};
\end{scope}
% --- analyst and the origin of Q2 ---
\path (expl.north |- server.north) ++(0,8.5mm)
  node[draw=black!75, line width=0.5pt, rounded corners=1.5pt, inner xsep=4pt, inner ysep=2.5pt]
  (analyst) {\textbf{Analyst}};
\node[draw=black!45, dashed, rounded corners=1.5pt, inner sep=2pt, text=black!65, align=center,
      anchor=west] (origin) at (server.west |- analyst) {refinement\\ mechanism};
\draw[flow, dashed, draw=black!55] (origin) -- node[above, inner sep=1.5pt, text=black!65] {$Q_2$} (analyst);
% --- what crosses the boundary of the server ---
\draw[flow] (analyst.south) -- (analyst.south |- server.north)
  node[midway, left, inner sep=2pt] {$Q_1,\,Q_2,\,\views,\,\Attrs,\,k,\,\epsilon$};
\draw[flow, draw=peeqorange, line width=0.7pt] (rel.north) |- (analyst.east)
  node[pos=0.75, above, inner sep=1.5pt, text=peeqorange] {$k$ explanations};
\end{tikzpicture}%

\caption{The \sys{} framework. The analyst supplies both queries, however $Q_2$ was obtained, with the search space and the budget. The server privately selects views, then explanations, and releases only noisy histogram pairs. The parameters $\tau$ and $\gamma$ are derived from public quantities and consume no budget.}
% (addressed: A_R added to the input, here and in figures/framework_tikz.tex) \yuval{Should $\mathcal{A}_R$ be part of the input as well?}
\label{fig:framework}
\end{figure}

\paratitle{The structure of an explanation.}
To say where the outputs differ and in what respect, an explanation names a region of the data and an aspect of it.
The region is given by a view predicate $p\in\views$, and the aspect is an attribute $A\in\Attrs$ whose distribution is compared between the two outputs inside the region.

\begin{definition}[Explanation candidate]
\label{def:candidate}
A \emph{candidate} is a pair $c=(p,A)$ of a view predicate $p\in\views$ and an explanation attribute $A\in\Attrs$ that $p$ does not constrain.
It induces the two histograms
\[
    \hist{Q_1(\pD)}{A}
    \qquad\text{and}\qquad
    \hist{Q_2(\pD)}{A}.
\]
Writing $\mathrm{attr}(p)$ for the set of attributes that $p$ constrains, the space of candidates is
\[
\mathcal{E}=\{(p,A)\in\views\times\Attrs \;:\; A\notin\mathrm{attr}(p)\},
\]
and we write $A_c$ for the attribute of a candidate $c\in\mathcal{E}$.
\end{definition}

We exclude the attributes that the view constrains because a view that fixes $A$ leaves both histograms on a single bin, and such a pair reports a change in one count rather than in a distribution.

\begin{example}[Explanation candidates]
\label{ex:candidate}
Consider the records of Figure~\ref{fig:intro-records} and the queries of Example~\ref{ex:intro}, where $Q_1$ selects tuples $1$, $2$ and $5$, and $Q_2$ selects tuples $3$, $4$ and $5$.
Let $c=(p,\texttt{occupation})$ with $p=(\texttt{education}=\texttt{Masters})$.
The histogram of $Q_1$ inside the view holds one executive (tuple $5$), and that of $Q_2$ holds one professional, one clerical worker and one executive (tuples $3$, $4$ and $5$).
Both outputs thus reach the view, and they differ in their occupational mix.
Now let $c'=(p',\texttt{occupation})$ with $p'=(\texttt{education}=\texttt{Bachelors})$.
The histogram of $Q_1$ holds one executive and one professional (tuples $1$ and $2$), and that of $Q_2$ is empty, since $Q_2$ selects only Master's degree holders.
We call such a view \emph{one-sided}.

Figure~\ref{fig:candidates} shows the same two candidates on the full dataset.
They illustrate two kinds of change.
Panel~(a) shows a shift inside a region that both outputs reach, and panel~(b) shows a region that disappears from the output as a whole.
\end{example}

What the analyst receives for a selected candidate is a noisy version of its two histograms.

\begin{definition}[Explanation]
\label{def:noisy-histogram}
For a set of tuples $T\subseteq D$, an attribute $A$ and a scale $b>0$, the \emph{noisy histogram} of $A$ over $T$ is
\[
\hat H_{T,A}(v)=\hist{T}{A}(v)+Y_v,\qquad v\in\Dom(A),
\]
where the variables $Y_v$ are drawn independently from the Laplace distribution of scale $b$, one for each bin of the public domain of $A$ (Definition~\ref{def:laplace-mechanism}).
The \emph{explanation} of a candidate $c=(p,A)$ is the pair $\big(\hat H_{Q_1(\pD),A},\,\hat H_{Q_2(\pD),A}\big)$.
\end{definition}

\paratitle{Problem statement.}
Given $D$, $Q_1$, $Q_2$, $\views$, $\Attrs$, $k$ and $\epsilon$, our goal is to release $k$ explanations of high quality while satisfying $\epsilon$-differential privacy.
We define the quality of an explanation next.


\subsection{Measures of Explanation Quality}
\label{sec:explanation-quality}

The running example has roughly ten thousand candidates, and the analyst can inspect only a few of them.\yuval{add the most important reference for each one here?}
Following the literature on explanations, we ask that the two histograms differ substantially (\emph{interestingness}), that the view be typical of the database (\emph{representativeness}) and hold a substantial part of the two outputs (\emph{coverage}), and that the $k$ explanations not repeat one another (\emph{diversity}).
The classical forms of these criteria fail in our setting for two reasons.
First, they are undefined or zero on one-sided views.
Second, they are ratios over groups that may be small, so a single tuple can change them substantially, and a private mechanism must add noise in proportion (Definition~\ref{def:global-sensitivity}).
For each measure, we show the failure on the running example and then our form.

\paratitle{Interestingness.}
Interestingness is commonly measured as a deviation~\cite{SilberschatzT96,GengHamilton06,TanKS02,Sarawagi98}, and the classical score of a candidate is the total variation distance $\TVD(P_1, P_2) = \frac{1}{2}\sum_{v} |P_1(v) - P_2(v)|$ between $P_i = P_{Q_i(p(D)),A}$, the distributions of $A$ over the two outputs inside the view~\cite{SeeDB,FEDEX}.
%$\TVD(P_1,P_2)=\frac12\sum_{x}|P_1(x)-P_2(x)|$ between the distributions of $A$ over the two outputs inside the view~\cite{SeeDB,FEDEX}, where $P_{T,A}(v)=\hist{T}{A}(v)/|T|$.
% On Figure~\ref{fig:intro-records}, it fails in both ways.
Both failures already appear on the records of Figure~\ref{fig:intro-records}.
For the one-sided candidate $c'$ of Example~\ref{ex:candidate} it is undefined, since $Q_2(p'(D))$ is empty.
For $c$ it is $2/3$, and it drops to $1/4$ when one professional with a Master's degree whom both queries select joins the database, since it normalizes by the output sizes inside the view.
We instead compare the two histograms in \emph{counts}, normalize each bin by the number of tuples that the view holds in it, which does not depend on the queries, and add a public constant $\gamma$ to every denominator.

\begin{definition}[Smoothed histogram distance]
\label{def:hd}
For a candidate $(p,A)$ and a public smoothing constant $\gamma \ge 0$,
\begin{multline}
\HD_\gamma(D;p,A,Q_1,Q_2) = \\
\sum_{v\in \Dom(A)}
\frac{\big|\hist{\Qone(\pD)}{A}(v)-\hist{\Qtwo(\pD)}{A}(v)\big|}
     {\hist{\pD}{A}(v)+\gamma} ,
\label{eq:hd-definition}
\end{multline}
where for $\gamma=0$ the sum is restricted to the bins with $\hist{\pD}{A}(v)>0$.
We write $\HD_\gamma(c)$ for short.
\end{definition}

Intuitively, for $\gamma=0$ a bin contributes the fraction of its tuples on which the two outputs disagree, so the score is defined on one-sided views as well.
It is $1+1+0=2$ for $c$, where the outputs agree only on the executive, and $1+1=2$ for $c'$.
The constant $\gamma$ acts as $\gamma$ additional tuples on which the outputs agree.
A bin with $n$ tuples thus keeps a fraction $n/(n+\gamma)$ of its contribution, which damps the sparse bins that a single tuple can swing and leaves well-populated bins almost unaffected.
Section~\ref{subsec:explanation-generation} sets $\gamma$ from public quantities.
Note that $\HD_\gamma$ measures change \emph{relative} to the population of the view, so the $1{,}579$ Master's degree holders that the refinement adds ($59\%$ of those in the data) weigh more than the $2{,}181$ Bachelor's degree holders that it removes ($27\%$).
The absolute size of a change enters through coverage.

\paratitle{Representativeness.}
A difference inside an unusual subset of the data may reflect the particular people in that subset rather than the change of the query, so we prefer views that are \emph{typical} of the database, the opposite of the unusualness that subgroup discovery rewards~\cite{Wrobel97,HerreraCGJ11}.
The natural score for a reference attribute $A$ is $1-\TVD\big(P_{\pD,A},\,P_{D,A}\big)$, and it has the same two defects.
It is undefined for an empty view, and it is unstable on a small one. In Figure~\ref{fig:intro-records},
% In Table~\ref{tab:intro-records}, 
the view $\texttt{hoursPerWeek}\ge 60$, which holds a single man, scores $0.6$ with \texttt{gender} as the reference attribute, and its score jumps to $1$ when one woman who works $60$ hours a week joins the database.
We apply the same remedy.

\begin{definition}[Smoothed representativeness]
\label{def:rep}
Let $c^A_v$ and $d^A_v$ be the number of tuples with $A$-value $v$ in $D$ and in $\pD$, respectively, let $N=|D|$ and $m=|\pD|$, and let $\gamma_R\ge0$ be a smoothing constant.
For a single attribute $A$,
\begin{equation}
\Rep^A_{\gamma_R}(p,D)=1-\frac{1}{2}\sum_{v\in\Dom(A)}\left|\frac{c^A_v}{N}-\frac{d^A_v}{m+\gamma_R}\right| ,
\label{eq:representativeness-proxy}
\end{equation}
and $\Rep_{\gamma_R}(p,D)$ is the average of $\Rep^A_{\gamma_R}(p,D)$ over $A\in\mathcal{A}_R(p)=\mathcal{A}_R\setminus\mathrm{attr}(p)$, the reference attributes that $p$ does not constrain.
\end{definition}

Intuitively, smoothing scales the profile of $A$ in the view by its \emph{fidelity} $\theta=m/(m+\gamma_R)$.
A large view is scored by its actual profile, whereas a small view scores close to the neutral value $\tfrac12$, so a handful of tuples can neither earn nor lose much.
With $\gamma_R=1$, for instance, the view above scores $0.75$ and then $0.83$, instead of $0.6$ and $1$.
Averaging over several reference attributes requires a view to be typical in several respects.
We assume $|\mathcal{A}_R|>\ell$, so that $\mathcal{A}_R(p)$ is never empty.
A larger $\gamma_R$ stabilizes the score but lowers the fidelity of every view (see \ifpaper the full version~\cite{fullversion}\else \Cref{app:proof-rep-range}\fi), so we set
\begin{equation}
\gamma_R=\alpha N,\qquad \alpha=\tfrac{1}{100},
\label{eq:gammaR-rule}
\end{equation}
and a view is read at fidelity $\tfrac12$ once it holds a fraction $\alpha$ of the database.
% \ag{Alon: no experiment varies $\alpha$ (the default $\gamma_R=N/100$ is used throughout Section~\ref{sec:experiments}). Please add a small sensitivity experiment, e.g.\ $\alpha\in\{1/1000,1/100,1/10\}$, under ``Other parameters'' in Section~\ref{subsec:quality}, and reference it here in one sentence, as we do for $\tau$ and for the budget split in Section~\ref{subsec:explanation-generation}.}
Unlike $\alpha$, the constant $\gamma_R$ is not public, because $N=|D|$ differs between neighboring databases, and Section~\ref{subsec:dp-proof} accounts for this.

\paratitle{Coverage.}
A difference on a handful of tuples is weak evidence, so explanation frameworks pair deviation with a measure of mass, such as support~\cite{AgrawalIS93,MacroBase,DIFF,Anchors,ExplanationTables}.
For a comparison, the natural choice is the two-sided support $|Q_1(\pD)|\cdot|Q_2(\pD)|$, which is $0$ for the one-sided Bachelor's view and which a single tuple changes by an amount that grows with the outputs.
We instead \emph{add} the two output sizes:
\begin{equation}
\Cov(p,Q_1,Q_2,D)=\frac{|Q_1(\pD)|+|Q_2(\pD)|}{|D|}.
\label{eq:coverage-proxy}
\end{equation}
Intuitively, coverage is the fraction of the database that the two outputs place in the view, counted once for each output, so it ranges in $[0,2]$ and a one-sided view keeps the full weight of its output.
For example, the Bachelor's view has coverage $2{,}181/48{,}842\approx0.045$, and the Master's view, with $874+2{,}453$ output tuples, has $0.068$.

\paratitle{Diversity.}
Explanations that are shown together should not repeat one another~\cite{MMR,GollapudiS09,LinBilmes11,DrosouP10}.
The aspect that an explanation describes is its attribute, so we discount an attribute by the number of times that it was already used:
\begin{equation}
\DivS(A,S)=\lambda^{\,\mathrm{count}(A,S)},
\qquad 0\le\lambda\le1,
\label{eq:diversity-score}
\end{equation}
where $S$ is the set of explanations selected so far and $\mathrm{count}(A,S)$ is the number of them with attribute~$A$.
The discount reads only released explanations and never the database, so it has no sensitivity of its own.

\paratitle{The quality of a set of explanations.}
The discount depends on the order of selection, whereas what we release, and wish to evaluate, is a set.

\begin{definition}[Set-level utility]
\label{def:setobj}
Let $E\subseteq\mathcal{E}$ be a set of candidates, and for an attribute $A$ let
$h^A_{(1)}\ge h^A_{(2)}\ge\cdots$ be the scores $\HD_\gamma(c)$ of the candidates $c\in E$ with $A_c=A$, in decreasing order, with $h^A_{(j)}=0$ for every $j$ beyond the number of such candidates.
For a non-increasing sequence of weights $w_1\ge w_2\ge\cdots\ge0$ with $w_1=1$,
\begin{equation}
F_w(E)=\sum_{A\in\Attrs}\;\sum_{j\ge1} w_j\, h^A_{(j)} .
\label{eq:set-objective}
\end{equation}
We write $F_\lambda$ for $F_w$ with the weights $w_j=\lambda^{\,j-1}$.
\end{definition}

Intuitively, among the explanations that share an attribute, the strongest counts in full and each further one is discounted by its rank.
With $\lambda=\tfrac12$, for example, two explanations on \texttt{occupation} of scores $9$ and $4$ and one on \texttt{age} of score $7$ have utility $9+\tfrac12\cdot4+7=18$.
$F_\lambda$ equals the total that the discount of Eq.~\eqref{eq:diversity-score} awards $E$ under the best ordering, and bounds it under every ordering (see \ifpaper the full version~\cite{fullversion}\else \Cref{app:proof-setobj}\fi).

\paratitle{Why the classical forms cannot be used.}
The failures above are not artifacts of a small table, as we show next.
Recall that $\Delta_f(N)$ denotes the largest change of a score $f$ between a database of size $N$ and a neighboring one (Section~\ref{subsec:prelim-dp}).

\begin{proposition}[Classical forms are sensitive]
\label{prop:sensitivity-itemized}
Over the neighboring databases on which the following scores are defined, (1)~the total variation distance $f$ between the distributions of $A$ over $Q_1(\pD)$ and over $Q_2(\pD)$ satisfies $\Delta_f(N)\ge\frac12$, which is half of its range, (2)~representativeness without smoothing ($\gamma_R=0$) satisfies $\Delta_f(N)\ge\frac12-\frac1N$, and (3)~the two-sided support $f=|Q_1(\pD)|\cdot|Q_2(\pD)|$ satisfies $\Delta_f(N)=2N+1$.
Hence none of the three sensitivities decreases as the database grows.
\end{proposition}

The proof appears in \ifpaper the full version~\cite{fullversion}\else \Cref{app:proof-classical}\fi.
In contrast, Section~\ref{subsec:dp-proof} shows that $\HD_\gamma$ has sensitivity $1/(1+\gamma)$, and that smoothed representativeness and coverage, combined into one score of a view, have sensitivity at most $4/N$.
